% Owner: Duy. Source evidence is the inspected backend snapshot; Listing and
% Booking is absent from its parent build and Compose file.
\section{Event-Driven Design and Infrastructure}\label{sec:events}

\subsection{Cross-Service Rental Workflow}
The specified rental lifecycle spans Listing and Booking, Contract and Tenant, Property, Billing and Utility, and User and Notification. Contract source contains a consumer for \texttt{RentalRequestAccepted}, but the corresponding Listing and Booking publisher is not integrated in the inspected backend tree, so the complete path is not source-verified. Each proposed transition has one authoritative owner; downstream services would update their own records rather than the publisher's database, consistent with local-transaction coordination \cite{richardson2018microservices}. The earlier workflow diagram is omitted because it labels unverified paths as verified.

In the target workflow, a rental request advances from \texttt{PENDING} to \texttt{VERIFIED} and \texttt{ACCEPTED}. Listing and Booking is intended to publish \texttt{RentalRequestAccepted} on \texttt{listing.events}; the inspected Contract consumer parses that event and invokes draft creation while saving a processed-event record in the same local transaction. Table~\ref{tab:workflow_states} separates the subsequent proposed transitions by owner and precondition. The source also retains \texttt{POST /internal/contracts/drafts}, so it cannot yet be assumed that Kafka is the only draft-creation path.

\begin{table*}[t]
\caption{Selected Target Workflow Transitions and Required Checks}
\label{tab:workflow_states}
\centering\footnotesize
\begin{tabularx}{\textwidth}{@{}p{0.18\textwidth}p{0.18\textwidth}p{0.19\textwidth}X@{}}
\toprule
Trigger and state owner & Precondition & Intended next state & Unresolved or exceptional outcome \\
\midrule
Request accepted; Contract & Accepted request; no prior draft for its identifier & One \texttt{DRAFT} contract & Duplicate event or synchronous draft call must not create a second draft; malformed events can enter a DLT. \\
Contract signed; Listing and Property & Contract becomes active & Listing closed; room \texttt{OCCUPIED} & Competing open requests and roommate posts require explicit resolution. \\
Signature declined or expired; Listing & Draft signature fails or expires & Request \texttt{VERIFIED} or rejected, respectively & Reopening a request must not be confused with reversing a signed contract. \\
Contract ended or move-out approved; Listing & No successor tenancy; cleaning rule satisfied & Listing reopened as \texttt{VACANT} or \texttt{UPCOMING\_VACANT} & Out-of-order or missing events leave the listing projection stale. \\
Tenant added or add rejected; Listing & Roommate application awaits settlement & Application accepted or rejected & Duplicate settlement must be idempotent. \\
\bottomrule
\end{tabularx}
\end{table*}

The specification assigns Billing a local contract summary and Property an occupancy projection. The inspected Billing and separate Property sources contain consumers for contract events, while User and Notification has a listener configured for the six named topics. These source observations do not show that the complete acceptance-to-billing sequence runs in Compose. Notification delivery also depends on optional e-mail or push configuration and is outside the draft-contract invariant. Contract source includes a scheduled meter-handover event path; its delivery frequency has not been measured.

\subsection{Synchronous Communication with OpenFeign}
The specification enumerates synchronous reads on \texttt{/internal/} endpoints for room capacity, tenant checks, billing data, and staff-building assignments. Contract, Billing, and Asset source contains Feign clients for some of these reads; Billing caches assignment results for five minutes. The inspected Contract module also exposes a mutating internal draft-creation endpoint. The gateway's path filter rejects external \texttt{/internal/} requests, but callers within the Docker network or on a published service port may reach internal APIs directly.

The common library specifies Feign timeouts, retry rules, error decoding, and header propagation for clients that use it. Eureka names are configured for service discovery. The design intends to avoid retrying mutating Feign commands, but the retained internal draft endpoint is an exception to a read-only interface rule. Source inspection indicates that Contract's occupant-add path requires a capacity read and Billing's invoice close depends on upstream data. Failure behavior for every internal call and circuit-breaker state still needs an endpoint-level test; these examples do not establish a universal fail-closed property.

\subsection{Asynchronous Communication with Apache Kafka}
The inspected Compose file configures one Kafka broker and an init container for six domain-event topics and matching dead-letter topics, each with three partitions and replication factor one. These are topics for the six specified publishing services; User and Notification is a consumer in this model. Topic creation does not show that every publisher is present or running: Listing and Booking is absent from the supplied parent build and Compose file. A domain event represents a local fact, while a consumer decides whether and how to update its own state.

The common-library envelope defines a UUID \texttt{eventId}, \texttt{eventType}, \texttt{version}, \texttt{occurredAt}, optional \texttt{correlationId}, and payload. The inspected Contract consumer ignores event types it does not handle, but rejects malformed envelopes and missing required payload fields. Compatibility therefore requires more than a version number: producers and consumers must agree on required fields, optional additions, unknown-field handling, and migration of renamed fields. The correlation identifier is intended to connect request and event logs; propagation, retention, and queryability across the complete workflow have not been demonstrated.

The inspected Contract Kafka configuration uses a service-name consumer group, disables auto-commit, and sets record acknowledgement. Keying an event by aggregate identifier can preserve order within one partition, but cannot order different topics or guarantee that a business effect commits. Redelivery is possible after failure. The acceptance consumer attempts to suppress duplicate draft creation by recording \texttt{eventId} with the business effect; a record sent to a dead-letter topic may remain unapplied until an explicit recovery procedure runs. The paper makes no system-wide exactly-once claim.

\subsection{Transactional Outbox and Idempotent Processing}
The outbox addresses the database-and-broker dual-write problem \cite{kleppmann2017designing}. Auth, Contract, Billing, Asset, and the separate Property source contain outbox-related code. The intended rule is to save publication intent in the same local transaction as the business change, then relay it to Kafka. This source presence does not establish that every specified publisher follows the rule: the Listing and Booking publisher is absent from the supplied backend tree, and end-to-end publication is untested. The earlier outbox diagram is omitted because it presents unverified redrive and idempotency behavior as unconditional.

The inspected Auth, Contract, and Billing modules contain scheduled relays; the separate Property source has a lease-based relay. In each intended path, a crash after broker acknowledgement but before marking the outbox row published can produce another broker record with the same event identifier. The Contract acceptance consumer writes its draft effect and \texttt{processed\_event(event\_id)} within one local transaction. A replay becomes a no-op only if the publisher reuses the same identifier, the first transaction committed, and every relevant side effect participates in that transaction or has its own deduplication rule. Concurrency and crash tests are needed to establish those conditions for the tested workflow.

The inspected Contract consumer configures three retries with one-second backoff and then routes an unhandled record to \texttt{listing.events.DLT}. Other consumers require separate verification. Dead-letter routing removes the failed record from ordinary processing; it does not repair the missing draft or prove safe redrive. The evaluation must count dead-lettered accepted requests separately, retain the original event identifier, and specify alerting and manual or automated redrive before claiming eventual completion. Scheduled domain jobs in the available modules are outside the acceptance-to-draft test.

\subsection{API Gateway and Request Processing}
Spring Cloud Gateway is the intended entry point for public business APIs and is published on port 8080. Explicit predicates map \texttt{/api/v1/} prefixes to \texttt{lb://} service names resolved through Eureka, and discovery-locator routing is disabled. Documentation routes are also configured. Other host-published ports are identified below; a gateway route does not demonstrate a reachable target.

The gateway source contains correlation-ID and JWT filters. It rejects external \texttt{/internal/} requests, removes client-supplied identity headers, whitelists selected public routes, and validates a Bearer token for other routes. Downstream authorization trusts the forwarded fields, making direct service reachability a security assumption discussed in Section~\ref{sec:architecture}. In the inspected Compose file, the gateway is published on host port 8080; User and Notification is additionally bound to host loopback port 8082. Kafka UI (8090), MinIO API (9000), and MinIO console (9001) are also host-published. These mappings are configuration facts, not evidence of HTTPS termination or successful access-control tests.

\subsection{Deployment and CI/CD Infrastructure}
The inspected \texttt{infra/docker-compose.yml} declares PostgreSQL with seven logical database names, one Kafka broker and topic-init container, MinIO with a report-images bucket init step, Eureka, the gateway, and three business containers: Auth, User and Notification, and Asset and Maintenance. Its bridge network connects the declared containers. Health checks and dependencies specify startup order but do not prove successful startup. The earlier deployment figure is omitted because it depicts eight business containers, an eighth database, four buckets, and a host-port boundary that the supplied Compose file does not configure.

The declared service Dockerfiles use Maven build stages and Java 21 runtime images. Compose reads credentials from environment variables; the checked-in files do not prove how deployed secrets are stored. A GitHub Actions workflow is configured for pushes to \texttt{main} or \texttt{master} and pull requests, with Java 21 and \texttt{mvn -B clean verify}. No successful run log was supplied for this snapshot, and the parent build omits Listing and Booking. Thus, CI configuration does not demonstrate a successful build of the complete rental workflow or production readiness.

% Owner: Duy. Logical event classes only; exact Kafka topic names are TODO.
\begin{table*}[t]
\caption{Candidate Event Topics and Publishers}
\label{tab:events}
\centering\footnotesize
\begin{tabularx}{\textwidth}{@{}p{0.19\textwidth}p{0.22\textwidth}X X@{}}
\toprule
Logical event or topic & Proposed publisher & State change represented & Contract details to verify \\
\midrule
Booking confirmed & Listing and Booking & Booking has passed the owner's confirmation rule & Topic name, key, schema, consumers \\
Contract activated & Contract and Tenant & Contract or tenancy has entered its active state & Legal trigger, topic name, key, consumers \\
Bill issued & Billing and Utility & A bill is available under the owner's rules & Billing trigger, topic name, schema, consumers \\
Maintenance requested & Asset and Maintenance & A maintenance request has been recorded & Topic name, key, schema, consumers \\
\bottomrule
\end{tabularx}
\end{table*}

